\section{Background and Related Work}
\label{sec:related}

\subsection{Vulnerability datasets for source code}
\label{sec:related:datasets}

The benchmarks the source-code vulnerability detection literature
trains and evaluates on are heavily skewed toward C and C++.
BigVul~\cite{TODO_bigvul} aggregates 264~CVEs across 348~open-source
C/C++ projects with paired vulnerable and fixed function pairs.
CVEfixes~\cite{TODO_cvefixes} extends this to multiple languages but
remains C-dominant in the resulting corpus, and the authors note that
Python and other interpreted languages contribute only a small
fraction of the labeled samples. CrossVul~\cite{TODO_crossvul} adds
broader language coverage (forty languages) but the per-language
sample counts for non-C languages are too small to support per-class
evaluation. DiverseVul~\cite{TODO_diversevul} explicitly deduplicates
across BigVul, Devign, ReVeal, CrossVul, and CVEfixes and reports that
roughly eighteen percent of records are duplicated when sources are
combined naively.

Two further datasets are worth naming. The Juliet test suite from
NIST~\cite{TODO_juliet} is hand-curated and synthetic, with one
vulnerable and one fixed version per CWE-pattern combination. It is
not a real-world distribution, and detectors that train on Juliet
typically transfer poorly to natural code~\cite{TODO_juliet_limits}.
Devign~\cite{TODO_devign} and ReVeal~\cite{TODO_reveal} are C-only
corpora drawn from Chromium, FFmpeg, Linux, and similar large
projects; they are influential for model-architecture comparisons
but tell us nothing about Python.

Within Python specifically, \texttt{vudenc}~\cite{TODO_vudenc} is the
closest existing labeled corpus. It contains human-annotated
vulnerable Python files across seven CWE classes (CWE-22, 78, 79, 89,
352, 601, 94). The vudenc data is high-quality at the file level but
does not carry paired \texttt{code\_after} versions, which restricts
the label-validation methodology that can be applied to it. We
include vudenc data in our benchmark as a separate label-confidence
tier (Section~\ref{sec:dataset:final}).

\subsection{Label noise in CVE-derived datasets}
\label{sec:related:labelnoise}

Mislabeled samples in CVE-fix-derived datasets are a recurring
problem in the literature. Croft et al.~\cite{TODO_croft} audited
BigVul and reported a label-noise rate of approximately twenty-five
percent on a manually reviewed sample. The root cause they identified
is the same one our audit surfaces: a CVE fix commit touches multiple
files, but only a subset of those files contains the actual
vulnerability. The remaining files are co-changed for refactoring,
test updates, documentation, or auxiliary cleanup, and inherit the
CVE's CWE label by accident.

Chakraborty et al.~\cite{TODO_reveal_critique} made a related
observation when evaluating ReVeal: they found that deep models
trained on commit-level labels learn to match surface patterns of the
co-changed files rather than the actual vulnerability pattern, and
that this is the dominant explanation for the poor cross-project
generalization that vulnerability detectors exhibit.

Neither study proposed an automated remediation. Our diff-based
filter (Section~\ref{sec:dataset:diff}) is, to our knowledge, the
first sink-anchored, evidence-based label-validity criterion in this
literature. The closest prior approach is D2A~\cite{TODO_d2a}, which
uses differential analysis between a buggy and a fixed version of a
program to assign labels, but D2A operates at the program-level rather
than the line-level and does not target sinks specifically.

\subsection{Static analysis tools for Python}
\label{sec:related:sast}

The static analysis landscape for Python is smaller than for C/C++
but includes several actively maintained tools. Bandit~\cite{TODO_bandit}
is the de-facto Python security linter; it ships approximately eighty
rules covering common CWE patterns and operates as a single-file
analyzer with no inter-procedural reasoning. Semgrep~\cite{TODO_semgrep}
is a multi-language pattern matcher with a YAML-based rule format and
a curated Python security ruleset (\texttt{p/python}, \texttt{p/security-audit}).
CodeQL~\cite{TODO_codeql} is GitHub's database-backed static analyzer;
it requires a built code database per project, which limits its
applicability at the file level, but it performs richer dataflow
analysis than the lighter-weight competitors. Pylint~\cite{TODO_pylint}
and Pyflakes are general-purpose linters that flag a small subset of
security-relevant patterns as a byproduct of code-quality checks.

All four tools fall into the broad category that
CASTLE~\cite{TODO_castle} labels ``Static Application Security
Tester'' or ``Generic Code Analyzer.'' They share a design philosophy:
detect by matching the source against a curated rule set. This makes
them fast and reproducible but exposes them to the surface-pattern
attacks that our four mutators (Section~\ref{sec:methodology})
specifically probe.

\subsection{Large language models for vulnerability detection}
\label{sec:related:llm}

Recent work uses large language models as vulnerability detectors
directly, prompting the model to classify or localize vulnerabilities
in source code. Khare et al.~\cite{TODO_khare_llm} evaluated GPT-3.5
and GPT-4 on C/C++ samples from existing benchmarks and reported
F1 scores competitive with fine-tuned transformer baselines, though
they noted high variance across prompt phrasings. Steenhoek
et al.~\cite{TODO_steenhoek} extended this comparison to ChatGPT,
Claude, and several open-source models, finding that LLM accuracy
scales with model size and that all models lose accuracy on
real-world (vs. synthetic) code.

The CASTLE benchmark~\cite{TODO_castle} is the most directly
relevant prior work. CASTLE evaluates fourteen static-analysis tools
and ten LLMs on a hand-curated 250-program C/C++ corpus with a
purpose-designed scoring formula (CASTLE Score) that combines true
positives, true negatives, false positives, and false negatives.
The headline CASTLE finding is that LLMs outperform every individual
SAST tool, with GPT-o3-Mini achieving the top score of 977 against a
theoretical maximum of 1{,}250. CASTLE does not measure adversarial
robustness, does not validate its labels beyond CVE inheritance, and
does not cover Python.

\subsection{Adversarial robustness of code models}
\label{sec:related:adv}

Code-model robustness against semantics-preserving perturbations is
an established subfield. Yefet et al.~\cite{TODO_yefet} showed that
identifier renaming alone is sufficient to flip the prediction of
state-of-the-art code-summarization and bug-detection models on a
substantial fraction of inputs. Bielik and Vechev~\cite{TODO_bielik}
formalized adversarial robustness for code as a refinement-type
problem and developed defenses based on input invariance. Henkel
et al.~\cite{TODO_henkel} applied program-transformation attacks
(variable renaming, statement insertion, code splitting) to neural
program-analysis tasks and reported large accuracy drops across the
board.

Our four mutators draw on this literature but are designed with a
different goal: rather than mounting the strongest possible attack
on a single model, we want a stable, interpretable benchmark for
comparing tools. Each mutator targets one specific shortcut
(positional, substring, identifier, single-function) so that a
per-mutator F1 drop is diagnostically meaningful. CASTLE evaluates
detection accuracy at one level; we add the perpendicular axis of
robustness under semantics-preserving rewriting.

\subsection{Positioning}
\label{sec:related:positioning}

Our work occupies the gap left by the above. We are not the first
to build a vulnerability dataset; BigVul, CVEfixes, and others have
done so at larger scale. We are not the first to identify label
noise in CVE-derived data; Croft et al.\ and the ReVeal authors
have. We are not the first to compare static analyzers against
LLMs; CASTLE did this for C/C++. We are not the first to study
robustness against code transformations; Yefet, Bielik, and Henkel
laid that groundwork.

What is new here is the combination. A Python benchmark with
seven sink-shaped Top-25 CWE classes does not currently exist.
An automated, sink-anchored label-validity criterion (the diff-based
filter) does not currently appear in the literature. And no prior
benchmark for vulnerability detection reports a per-mutator
robustness breakdown alongside its detection accuracy. We argue
that the combination is what makes the benchmark useful: every
component has been studied in isolation, but the integrated
evaluation framework is what surfaces the diagnostic information
about which detectors do real reasoning and which do pattern
matching.
